\section{\gai Acceleration}
\label{section:34-related-gai-acceleration}

elevate vs exploit



\textbf{Graph Pre-Processing}
Many of the computational problems are not from the graph processing itself, but from generating heuristics and optimizing approximate solutions. 
One option to mitigate the impact of poor spatial locality is to pre-process the graph and re-order the vertices in memory which can be computationally expensive and time-consuming~\cite{Balaji2018}
Another is to use vertex scheduling to determine the optimal order to visit vertices to maximize spatial locality. 
vertex scheduling can be offline or online, though online approaches are generally infeasible because of their computational cost~\cite{Mukkara2018}.
Another optimization that aims to reduce memory bandwidth consumption and latency is `cache-blocking' where the graph is decomposed to subgraphs that can fit into the available cache. 
Closely related is the idea of \textit{graph-partitioning} which seeks to optimize the division of the graph between worker nodes in a parallel setting so that there are minimal cross-worker dependencies~\cite{Balaji2018}.
As is discussed in section~\ref{section:124-general-characteristics}, optimal graph partitioning is $NP$-Complete.
But sub-optimal partitioning leads to load imbalance and contention between workers for high-degree vertices~\cite{Chen2019}.
So, heuristic approaches approximate the partitions.
These approaches can involve computationally expensive eigendecompositions~\cite{Ng2001}, or complex test-and-adjust approaches like the \textit{Kernighan-Lin} algorithm~\cite{Kernighan1970} among others.
The shape of the graph has a significant impact on the `hardness' of partitioning.
\textit{Scale-free} graphs with small diameters are very difficult to partition, whereas \textit{mesh} graphs with large diameters are more easily separable because of the reduced interdependence of nodes~\cite{Beamer2015}. 
Indeed, there is a reasonable assertion that the key to improving the effectiveness of graph processing algorithms may depend on finding algorithms that can efficiently partition graphs and process them as independent subgraphs including streaming and distributed partitioning techniques~\cite{Mccune2015}.


\subsubsection{Paralellism as a solution}
Beginning with distributed approaches, and growing to include shared-memory parallelism when single machines gained the memory capacity to support processing, parallel approaches to graph processing are essential to \textit{efficient} and \textit{effective} graph processing. 
It is particularly hampered by load imbalance, synchronization overheads and non-uniform memory accesses (NUMA)~\cite{Beamer2015}

\textbf{Data and Algorithm Dependence.}
Parallel algorithms are particularly sensitive to the shape of the input data. 
Highly skewed degree distributions with small diameters are difficult to partition effectively~\cite{Beamer2015, Wu2015}.
Algorithms that are \textit{compute intensive} like PageRank which need to process all of the vertices in each iteration suffer from poor locality, and the memory latency cannot be hidden because of the poor communication-computation ratio~\cite{Ahn2015}.  

\textbf{Parallel abstractions.}
In general, graph computations are difficult to parallelize using traditional approaches because of their irregular structure and poor communication-computation ratio~\cite{Buluc2011}.
Factors like skewed degree distributions of graphs, heterogeneity of input graph types (scale-free, random, mesh) and maintaining synchronization make developing correct and efficient algorithms very difficult~\cite{Hong2012, Ediger2013}.
The proliferation of parallel graph processing frameworks and \textit{Domain Specific Languages} (DSL) that aim to abstract away all the hard-to-implement parallel processing operations and builds in expert optimizations so that the end-user-developer doesn't have to~\cite{Nguyen2013, Dhulipala2017}.
A common approach is to decompose graph operations into a collection of \textit{primitive} or \textit{atomic} operations, and then combine those to compose the required algorithms for processing. 
Vertex-centric models~\cite{Mccune2015}, edge-centric models, linear-algebraic models~\cite{Mattson2013, Kepner2016, Buluc2011, Azad2021}, models that apply the \textit{Gather Apply Scatter} (GAS), or the \textit{Frontier} model~\cite{Wang2016} all provide an abstraction to facilitate more efficient parallel graph processing. 
However, the frameworks are limited, with certain classes of algorithms unable to be represented in their abstractions (e.g. SSSP or MST implementations that use priority queues in early linear-algebraic approaches~\cite{Buluc2011}, bucketing algorithms in \textit{frontier} based models~\cite{Dhulipala2017}), impacts on message passing speed and memory consumption~\cite{Wu2015}, and a general finding that an abstracted library does not perform as \textit{efficently} as a hand-tuned natively implemented algorithm optimized for the system and data it is being run on~\cite{Satish2014}.

\textbf{Distributed vs. Shared-Memory}
The tradeoffs between distributed and shared memory parallelism have been discussed extensively above, but the core of the tradeoff is that distributed systems incur a network bottleneck that does not impact shared-memory systems in return for accessing potentially greater memory and compute resources~\cite{Shun2013, Beamer2015}.
Shared memory systems can exacerbate the already limited memory bandwidth issues as more compute resources compete for the limited memory bandwidth~\cite{Shun2015}, and so does not scale with the number of additional compute nodes~\cite{Ahn2015,Kepner2016}.

\textbf{Data parallelism vs Amorphous Data Paralellism}.
Most general-purpose parallel computing frameworks were developed for data parallelism. 
Data parallelism allows the same operation to be performed on multiple data at the same time to increase efficiency. 
While some graph processing algorithms have data parallelism, many exhibit a more complex parallelism pattern called “amorphous” data-parallelism~\cite{Lenharth2016}.
\textit{Amorphous} data parallelism is a more flexible approach that can better handle irregularity in the shape of data, at the cost of being processed less efficiently by general-purpose hardware. 

\subsubsection{GPU Acceleration}
The use of GPUs to accelerate computation has seen enormous success in the deep learning community in recent years~\cite{Dally2021}.
GPUs have significantly more memory bandwidth than a standard CPU architecture and are optimized for dense, regular computations. 
The increased memory bandwidth allows them to \textit{brute-force} their way through some of the memory-boundedness of CPU-based implementations of algorithms~\cite{Aananthakrishnan2020}, but they have several notable drawbacks. 

\textbf{GPU Memory.}
GPUs use high-speed on-chip memory to accelerate their operations. 
As a result, the data that they are processing needs to fit in the on-chip memory, which is typically much smaller than the main memory of the host system. 
The limited memory capacity makes it infeasible for them to process large graphs as a single node~\cite{Aananthakrishnan2020}, meaning that they need to be used as part of a parallel processing framework and manage the associated baggage for parallel processing systems discussed above.
The limited memory availability makes it particularly difficult to implement algorithms with large intermediate representations, as they consume the available memory and fail~\cite{Sundaram2015}.
The memory issue is exacerbated by the `GPU Trick' of densifying the connectivity matrix for more efficient computation which can increase memory consumption by orders of magnitude. 
The sparsity of graphs can also interfere with the memory coalescing process that is critical to efficient GPU computation~\cite{Aananthakrishnan2020}.
As a result, GPU graph programming frameworks focus on how to manage the number of `active' vertices during a computation to reduce the demands on the memory~\cite{Wang2016}.
It also highlights the importance of efficient sparse representations, with the CSR representation effectively doubling the size of the graph that can be processed by the GPU compared to a sparse adjacency matrix representation~\cite{Liu2014}, which is why it is the favored data format for all BSP implemenations~\cite{Wu2015}.

It is also important to note that the host system must copy the data across to the GPU before it can be processed, so for the use of an accelerator to be beneficial the efficiency gain in execution time must outweigh the cost of the partitioning and data movement~\cite{Azad2021}.

\textbf{GPU Compute.}
GPU computing is optimized for dense-regular computations that heavily exploit data parallelism. 
Graph processing involves sparse-irregular computations that feature amorphous data parallelism.
The impact of input data shape is evident, with more regular degree distributions (like a mesh network) seeing throughput increase linearly as the mean degree increases, but falling sharply as dense, skewed distributions inflict load-balancing overheads~\cite{Wu2015}.
As a result, GPUs are not optimized for graph processing, struggling with irregular data accesses/control flow~\cite{Wang2016}.
The Bulk Synchronous Parallel (BSP) paradigm employed by most GPU programming frameworks correlates the number of atomic (primitive) operation invocations with the overall performance of the algorithm--approaches that minimize the work they do tend to perform better overall~\cite{Wu2015}.
The primary challenge when using GPUs to execute graph algorithms is workload distribution~\cite{Wang2016}.
The irregular processing patterns mean that the thread assignments for the GPU can't be known until execution time because of their dependence on the data. 
Thread assignment is a common step in GPU pre-processing and deferring to execution time degrades GPU performance relative to a typical dense, regular workload~\cite{Fu2014}.
Imbalance in the graph because of a skewed degree distribution exacerbates the unpredictability of load-balancing that comes from not being able to pre-assign the threads and can result in uncoalesced memory patterns, and thread divergence, severely hampering GPU performance relative to dense regular tasks~\cite{Fu2014}. 
In a distributed setting, branches cause thread divergence and synchronization overheads limit thread progress~\cite{Aananthakrishnan2020}.

